\chaptertitle{Lisp: Atoms and Lists}

                          Lisp: Atoms and Lists
                                    February, 1983


IN previous columns I have written quite often about the field of artificial
intelligence-the search for ways to program computers so that they might come to
behave with flexibility, common sense, insight, creativity, self awareness, humor, and
so on. The quest for AI started in earnest over two decades ago, and since then has
bifurcated many times, so that today it is a very active and multifaceted research area.
In the United States there are perhaps a couple of thousand people professionally
involved in AI, and there are a similar number abroad. Although there is among these
workers a considerable divergence of opinion concerning the best route to Al, one
thing that is nearly unanimous is the choice of programming language. Most AI
research efforts are carried out in a language called "Lisp". (The name is not quite an
acronym; it stands for "list processing".)
        Why is most AI work done in Lisp? There are many reasons, most of which
are somewhat technical, but one of the best is quite simple: Lisp is crisp. Or as
Marilyn Monroe said in The Seven-Year Itch, "I think it's jus-teeegant!" Every
computer language has arbitrary features, and most languages are in fact overloaded
with them. A few, however, such as Lisp and Algol, are built around a kernel that
seems as natural as a branch of mathematics. The kernel of Lisp has a crystalline
purity that not only appeals to the esthetic sense, but also makes Lisp a far more
flexible language than most others. Because of Lisp's beauty and centrality in this
important area of modern science, then, I have decided to devote a trio of columns to
some of the basic ideas of Lisp.
        The deep roots of Lisp lie principally in mathematical logic. Mathematical
pioneers such as Thoralf Skolem, Kurt Godel, and Alonzo Church contributed
seminal ideas to logic in the 1920's and 1930's that were incorporated decades later
into Lisp. Computer programming in earnest began in the 1940's, but so-called
"higher-level" programming languages (of which Lisp is one) came into existence
only in the 1950's. The earliest list-processing language was not Lisp but IPL
("Information Processing Language"), developed in the mid-1950's by Herbert Simon,
Allen Newell, and J. C. Shaw. In the years 1956-58, John McCarthy, drawing on all
these




Lisp: Atoms and Lists                                                               396

previous sources, came up with an elegant algebraic list-processing language he called
Lisp. It caught on quickly with the young crowd around him at the newly-formed MIT
Artificial Intelligence Project, was implemented on the IBM 704, spread to other AI
groups, infected them, and . has stayed around all these years. Many dialects now
exist, but all of them share that central elegant kernel.

                                     *       *      *

        Let us now move on to the way Lisp really works. One of the most appealing
features of Lisp is that it is interactive, as contrasted with most other higher-level
languages, which are noninteractive. What this means is the following. When you
want to program in Lisp, you sit down at a terminal connected to a computer and you
type the word "lisp" (or words to that effect). The next thing you will see on your
screen is a so-called "prompt" -a characteristic symbol such as an arrow or asterisk. I
like to think of this prompt as a greeting spoken by a special "Lisp genie", bowing
low and saying to you, "Your wish is my command-and now, what is your next
wish?" The genie then waits for you to type something to it. This genie is usually
referred to as the Lisp interpreter, and it will do anything you want but you have to
take great care in expressing your desires precisely, otherwise you may reap some
disastrous effects. Shown below is the prompt, then gn that the Lisp genie is ready to
do your bidding:

:>

The genie is asking us for our heart's desire, so let us type in a simple expression:

:> (plus 2 2)

and then a carriage return. (By the way, all Lisp expressions and words will be printed
in Helvetica in this and the following two chapters.) Even non-Lispers can probably
anticipate that the Lisp genie will print in return. the value 4. Then it will also print a
fresh prompt, so that the screen will now appear this way:

:> (plus 2 2)
:>

        The genie is now ready to carry out our next command-or, more politely
stated, our next wish-should we have one. The carrying-out of a wish expressed as a
Lisp statement is called evaluation of that statement. The preceding short interchange
between human and computer exemplifies the




Lisp: Atoms and Lists                                                                   397

behavior of the Lisp interpreter: it reads a statement, evaluates it, prints the
appropriate value, and then signals its readiness to read a new statement. For this
reason, the central activity of the Lisp interpreter is referred to as the read-eval-print
loop.
        The existence of this Lisp genie (the Lisp interpreter) is what makes Lisp
interactive. You get immediate feedback as soon as you have typed a "wish" -a
complete statement-to Lisp. And the way to get a bunch of wishes carried out is to
type one, then ask the genie to carry it out, then type another, ask the genie again, and
so on.
        By contrast, in many higher-level computer languages you must write out an
entire program consisting of a vast number of wishes to be carried out in some
specified order. What's worse is that later wishes usually depend strongly on the
consequences of earlier wishes-and of course, you don't get to try them out one by
one. The execution of such a program may, needless to say, lead to many unexpected
results, because so many wishes have to mesh perfectly together. If you've made the
slightest conceptual error in designing your wish list, then a total foul-up is likely-in
fact, almost inevitable. Running a program of this sort is like launching a new space
probe, untested: you can't possibly have anticipated all the things that might go
wrong, and so all you can do is sit back and watch, hoping that it will work. If it fails,
you go back and correct the one thing the failure revealed, and then try another
launch. Such a gawky, indirect, expensive way of programming is in marked contrast
to the direct, interactive, one-wish-at-atime style of Lisp, which allows "incremental"
program development and debugging. This is another major reason for the popularity
of Lisp.

                                        *   *   *

        What sorts of wishes can you type to the Lisp genie for evaluation, and what
sorts of things will it print back to you? Well, to begin with, you can type arithmetical
expressions expressed in a rather strange way, such as (times (plus 6 3) (difference 6
3)). The answer to this is 27, since (plus 6 3) evaluates to 9, and (difference 6 3)
evaluates to 3, and their product is 27. This notation, in which each operation is
placed to the left of its operands, was invented by the Polish logician Jan >
ukasiewicz before computers existed. Unfortunately for > ukasiewicz, his name was
too formidable-looking for most speakers of English, and so this type of notation
came to be called Polish notation. Here is a simple problem in this notation for you, in
which you are to play the part of the Lisp genie:

:> (quotient (plus 2113) (difference 23 (times 2 (difference 7 (plus 2 2)))))

       Perhaps you have noticed that statements of Lisp involve parentheses. A
profusion of parentheses is one of the hallmarks of Lisp. It is not uncommon to see an
expression that terminates in a dozen right parentheses! This




Lisp: Atoms and Lists                                                                 398

makes many people shudder at first-and yet once you get used to their characteristic
appearance, Lisp expressions become remarkably intuitive, even, charming, to the
eye, especially when pretty printed, which means that a careful indentation scheme is
followed that reveals their logical structure. All of the expressions in displays in this
article have been pretty-printed.
        The heart of Lisp is its manipulable structures. All programs in Lisp work by
creating, modifying, and destroying structures. Structures come in two types: atomic
and composite, or, as they are usually called, atoms and lists. Thus, every Lisp object
is either an atom or a list (but not both). The only exception is the special object called
nil, which is both an atom and a list. More about nil in a moment. What are some
other typical Lisp atoms? Here are a few:

                    hydrogen, helium, j-s-bach, 1729, 3.14159, pi,
                         arf, foo, bar, baz, buttons-\&-bows

        Lists are the flexible data structures of Lisp. A list is pretty much what it
sounds like: a collection of some parts in a specific order. The parts of a list are
usually called its elements or members. What can these members be? Well, not
surprisingly, lists can have atoms as members. But just as easily, lists can contain lists
as members, and those lists can in turn contain other lists as members, and so on,
recursively. Oops! I jumped the gun with that word. But no harm done. You certainly
understood what I meant, and it will prepare you for a more technical definition of the
term to come later.
        A list printed on your screen is recognizable by its parentheses. In Lisp,
anything bounded by matching parentheses constitutes a list. So, for instance, (zonk
blee strill (croak flonk)) is a four-element list whose last element is itself a two-
element list. Another short list is (plus 2 2), illustrating the fact that Lisp statements
themselves are lists. This is important because it means that the Lisp genie, by
manipulating lists and atoms, can actually construct new wishes by itself. Thus the
object of a wish can be the construction-and subsequent evaluation-of a new wish!
        Then there is the empty list-the list with no elements at all. How is this written
\_ down? You might think that an empty pair of parentheses-()would work. Indeed, it
will work-but there is a second way of indicating the empty list, and that is by writing
nil. The two notations are synonymous, although nil is more commonly written than
() is. The empty list, nil, is a key concept of Lisp; in the universe of lists, it is what
zero is in the universe of-numbers. To use another metaphor for nil, it is like the earth
in which all structures are rooted. But for you to understand what this means, you will
have to wait a bit.
        The most commonly exploited feature of an atom is that it has (or can be
given) a value. Some atoms have permanent values, while others are variables. As you
might expect, the value of the atom 1729 is the integer




Lisp: Atoms and Lists                                                                  399

1729, and this is permanent. (I am distinguishing here between the atom whose print
name or pname is the four-digit string 1729, and the eternal Platonic essence that
happens to be the sum of two cubes in two different ways-i.e., the number 1729.) The
value of nil is also permanent, and it is -nil! Only one other atom has itself as its
permanent value, and that is the special atom t.
        Aside from t, nil, and atoms whose names are numerals, atoms are generally
variables, which means that you can assign values to them and later change their
values at will. How is this done? Well, to assign the value 4 to the atom pie, you can
type to the Lisp genie (setq pie 4). Or you could just as well type (setq pie (plus 2
2))-or even (setq pie (plus 1 1 1 1)). In any of these cases, as soon as you type your
carriage return, pie's value will become 4, and so it will remain forevermore-or at
least until you do another setq operation on the atom pie.
        Lisp would not be crisp if the only values atoms could have were numbers.
Fortunately, however, an atom's value can be set to any kind of Lisp object -any atom
or list whatsoever. For instance, we might want to make the value of the atom pi be a
list such as (a b c) or perhaps (plus 2 2) instead of the number 4. To do the latter, we
again use the setq operation. To illustrate, here follows a brief conversation with the
genie:

       :> (setq pie (plus 2 2)) 4
       :> (setq pi '(plus 2 2)) (plus 2 2)

         Notice the vast difference between the values assigned to the atoms pie and pi
as a result of these two wishes asked of the Lisp genie, which differ merely in the
presence or absence of a small but critical quote mark in front of the inner list (plus 2
2). In the first wish, containing no quote mark, that inner (plus 2 2) must be evaluated.
This returns 4, which is assigned to the variable pie as its new value. On the other
hand, in the second wish, since the quote mark is there, the list (plus 2 2) is never
executed as a command, but is treated merely as an inert lump of Lispstuff, much like
meat on a butcher's shelf. It is ever so close to being "alive", yet it is dead. So the
value of pi in this second case is the list (plus 2 2), a fragment of Lisp code. The
following interchange with the genie confirms the values of these atoms.

:> pie
:> pi
(plus 2 2)
:> (eval pi)
:>




Lisp: Atoms and Lists                                                                400

What is this last step? I wanted to show how you can ask the genie to evaluate the
value of an expression, rather than simply printing the value of that expression.
Ordinarily, the genie automatically performs just one level of evaluation, but by
writing eval, you can get a second stage of evaluation carried out. (And of course, by
using eval over and over again, you can carry this as far as you like.) This feature
often proves invaluable, but it is a little too advanced to discuss further at this stage.

         Every list but nil has at least one element. This first element is called the list's
Car. Thus the car of (eval pi) is the atom eval. The cars of the lists (plus 2 2), (setq x
17), (eval pi), and (car pi) are all names of operations, or, as they are more commonly
called in Lisp, functions. The car of a list need not be the name of a function; it need
not even be an atom. For instance, ((1)(2 2) (3 3 3)) is a perfectly fine list. Its car is
the list (1), whose car in turn is not a function name but merely a numeral.
         If you were to remove a list's car, what would remain? A shorter list. This -is
called the list's cdr, a word that sounds about halfway between "kidder" and
"could'er". (The words "car" and "cdr" are quaint relics from the first implementation
of Lisp on the IBM 704. The letters in "car" stand for "Contents of the Address part of
Register" and those in "cdr" for "Contents of the Decrement part of Register referring
to specific hardware features of that machine, now long since irrelevant.) The cdr of
(a b c d) is the list (b c d), whose cdr is (c d), whose cdr is (d), whose cdr is nil. And
nil has no cdr, just as it has no car. Attempting to take the car or cdr of nil causes (or
should cause) the Lisp genie to cough out an error message, just as attempting to
divide by zero should evoke an error message.
         Here is a little table showing the car and cdr of a few lists, just to make sure
the notions are unambiguous.

       list                    car                     cdr
       ((a) b (c))             (a)                     (b (c))
       (plus 2 2)              plus                    (22)
       ((car x) (car y))       (car x)                 ((car y))
       (nil nil nil nil)       nil                     (nil nil nil)
       (nil)                   nil                     nil
       nil                     **ERROR**               **ERROR**

        Just as car and cdr are called functions, so the things that they operate on are
called their arguments. Thus in the command (plus pie 2), plus is the function name,
and the arguments are the atoms pie and 2. In evaluating this command (and most
commands), the genie figures out the values of the arguments, and then applies the
function to those values. Thus, since the




Lisp: Atoms and Lists                                                                   401

value of the atom pie is 4, and the value of the atom 2 is 2, the genie returns the atom
6.

                                        *   *   *

        Suppose you have a list and you'd like to see a list just like it, only one
element longer. For instance, suppose the value of the atom x is (cake cookie) and
you'd like to create a new list called y just like x, except with an extra atom-say pie-at
the front. You can then use the function called cons (short for "construct"), whose
effect is to make a new list out of an old list and a suggested car. Here's a transcript of
such a process:

->(setq x '(cake cookie))
(cake cookie)
->(setq y (cons 'pie x))
(pie cake cookie)
-> x
(cake cookie)

        Two things are worth noticing here. I asked for the value of x to be printed out
after the cons operation, so you could see that x itself was not changed by the cons.
The cons operation created a new list and made that list be the value of y, but left x
entirely alone. The other noteworthy fact is that I used that quote mark again, in front
of the atom pie. What if I had not used it? Here's what would have happened.

-> (setq z (cons pie x))
(4 cake cookie)

Remember, after all, that the atom pie still has the value 4, and whenever the genie
sees an unquoted atom inside a wish, it will always use the value belonging to that
atom, rather than the atom's name. (Always? Well, almost always. I'll explain in a
moment. In the meantime, look for an exception - you've already encountered it.)
         Now here are a few exercises-some a bit tricky-for you. Watch out for the
quote marks! Oh, one last thing: I use the function reverse, which produces a list just
like its argument, only with its elements in reverse order. For instance, the genie, upon
being told (reverse '((a b) (c d e))) will write ((c d e) (a b)). The genie's lines in this
dialogue are given afterward.

-> (setq w (cons pie '(cdr z)))
-> (setq v (cons 'pie (cdr z)))
-> (setq u (reverse v))
-> (cdr (cdr u))
-> (car (cdr u))
-> (cons (car (cdr u)) u)
-> u




Lisp: Atoms and Lists                                                                  402

-> (reverse `(cons (car u) (reverse (cdr u))))
-> (reverse (cons (car u) (reverse (cdr u))))
-> u
-> (cons 'cookie (cons 'cake (cons 'pie nil)))

Answers (as printed by the genie):

(4 cdr z)
(pie cake cookie)
(cookie cake pie)
(pie)
cake
(cake cookie cake pie)
(cookie cake pie)
((reverse (cdr u)) (car u) cons)
(cake pie cookie)
(cookie cake pie)
(cookie cake pie)

         The last example, featuring repeated use of cons, is often called, in Lisp slang
"coning up a list". You start with nil, and then do repeated cons operations. It is
analogous to building a positive integer by starting at zero and hen performing the
successor operation over and over again. However, whereas at any stage in the latter
process there is a unique way of performing the successor operation, given any list
there are infinitely many different items you can cons onto it, thus giving rise to a vast
branching tree of lists instead of the unbranching number line. It is on account of this
image of a tree growing out of the ground of nil and containing all possible lists that I
earlier likened nil to "the earth in which all structures are rooted". ,
         As I mentioned a moment ago, the genie doesn't always replace (unquoted)
atoms by their values. There are cases where a function treats its arguments, though
unquoted, as if quoted. Did you go back and find such a case? It's easy. The answer is
the function setq. In particular, in a setq command, the first atom is taken straight-not
evaluated. As a matter of fact, the q in setq stands for "quote", meaning that the first
argument is treated as if quoted. Things can get quite tricky when you learn about set,
a function similar to setq except that it does evaluate its first argument. Thus, if the
value of the atom x is the atom k, then saying (set x 7) will not do anything to x-its
value will remain the atom k-but the value of the atom k will now become 7. So
watch closely:

-> (setq a 'b)
-> (setq b 'c)
-> (setq c 'a)
-> (set a c)
-> (set c b)




Lisp: Atoms and Lists                                                                 403

Now tell me: What are the values of the atoms a, b, and c? Here comes the answer, so
don't peek. They are, respectively: a, a, and a. This may seem a bit confusing. You
may be reassured to know that in Lisp, set is not very commonly used, and such
confusions do not arise that often.

                                        *   *   *

        Psychologically, one of the great powers of programming is the ability to
define new compound operations in terms of old ones, and to do this over and over
again, thus building up a vast repertoire of ever more complex operations. It is quite
reminiscent of evolution, in which ever more complex molecules evolve out of less
complex ones, in an ever-upward spiral of complexity and creativity. It is also quite
reminiscent of the industrial revolution, in which people used very simple early
machines to help them build more complex machines, then used those in turn to build
even more complex machines, and so on, once again in an ever-upward spiral of
complexity and creativity. At each stage, whether in evolution or revolution, the
products get more flexible and more intricate, more "intelligent" and yet more
vulnerable to delicate "bugs" or breakdowns.
        Likewise with programming in Lisp, only here the "molecules" or "machines"
are now Lisp functions defined in terms of previously known Lisp functions. Suppose,
for instance, that you wish to have a function that will always return the last element
of a list, just as car always returns the first element of a list. Lisp does not come
equipped with such a function, but you can easily create one. Do you see how? To get
the last element of a list called lyst, you simply do a reverse on lyst and then take the
car of that: (car (reverse lyst)). To dub this operation with the name rac (car
backwards), we use the def function, as follows:

-> (def rac (lambda (lyst) (car (reverse lyst))))

Using def this way creates a function definition. In it, the word lambda followed by
(lyst) indicates that the function we are defining has only one parameter, or dummy
variable, to be called lyst. (It could have been called anything; I just happen to like the
atom lyst.) In general, the list of parameters (dummy variables) must immediately
follow the word lambda. After this "def wish" has been carried out, the rac function
is as well understood by the genie as is car. Thus (rac '(your brains)) will yield the
atom brains. And we can use rac itself in definitions of yet further functions. The
whole thing snowballs rather miraculously, and you can quickly become
overwhelmed by the power you wield.
Here is a simple example. Suppose you have a situation where you know you are
going to run into many big long lists and you know it will often be useful to form, for
each such long list, a short list that contains just its car and rac. We can define a one-
parameter function to do this for you:




Lisp: Atoms and Lists                                                                  404

-> (def readers-digest-condensed-version
  (lambda (biglonglist)
   (cons (car biglonglist) (cons (rac biglonglist) nil))))

Thus if we apply our new function readers-digest-condensed-version to the entire
text of James Joyce's Finnegans Wake (treating it as a big long list of words), we will
obtain the shorter list (riverrun the). Unfortunately, reapplying the condensation
operator to this new list will not simplify it any further.
It would be nice as well as useful if we could create an inverse operation to readers-
digest-condensed-version called rejoyce that, given any two words, would create a
novel beginning and ending with them, respectively -and such that James Joyce would
have written it (had he thought of it). Thus execution of the Lisp statement (rejoyce
'Stately 'Yes) would result in the Lisp genie generating from scratch the entire novel
Ulysses. Writing this function is left as an exercise for the reader. To test your
program, see what it does with (rejoyce 'karma 'dharma).

                                       *   *   *

        One goal that has seemed to some people to be both desirable and feasible
using Lisp and related programming languages is (1) to make every single statement
return a value and (2) to have it be through this returned value and only through it that
the statement has any effect. The idea of (1) is that values are handed "upward" from
the innermost function calls to the outermost ones, until the full statement's value is
returned to you. The idea of (2) is that during all these calls, no atom has its value
changed at all (unless the atom is a dummy variable). In all dialects of Lisp known to
me, (1) is true, but (2) is not necessarily true.
        Thus if x is bound to (a b c d e) and you say (car (cdr (reverse x))), the first
thing that happens is that (reverse x) is calculated; then this value is handed "up" to
the cdr function, which calculates the cdr of that list; finally, this shorter list is
handed to the car function, which extracts one element-namely the atom d-and
returns it. In the meantime, the atom x has suffered no damage; it is still bound to (a b
c d e).
It might seem that an expression such as (reverse x) would change the value of x by
reversing it, just as carrying out the oral command "Turn your sweater inside out" will
affect the sweater. But actually, carrying out the wish (reverse x) no more changes
the value of x than carrying out the wish (plus 2 2) changes the value of 2. Instead,
executing (reverse x) causes a new (unnamed) list to come into being, just like x, only
reversed. And that list is the value of the statement; it is what the statement returns.
The value of x itself, however, is untouched. Similarly, evaluating (cons 5 pi) will not
change the list named pi in the slightest; it merely returns a new list with 5 as its car
and whatever pi's value is as its cdr.




Lisp: Atoms and Lists                                                                405

         Such behavior is to be contrasted with that of functions that leave "side
effects" in their wake. Such side effects are usually in the form of changed variable
bindings, although there are other possibilities, such as causing input or output to take
place. A typical "harmful" command is a setq, and proponents of the "applicative"
school of programming-the school that says you should never make any side effects
whatsoever-are profoundly disturbed by the mere mention of setq. For them, all
results must come about purely by the way that functions compute their values and
hand them to other functions.
         The only bindings that the advocates of the applicative style approve of are
transitory "lambda bindings"-those that arise when a function is applied to its
arguments. Whenever any function is called, that function's dummy variables
temporarily assume "lambda bindings". These bindings are just like those caused by a
setq, except that they are fleeting. That is, the moment the function is finished
computing, they go away-vanishing without a trace. For example, during the
computation of (rac '(a b c)), the lambda binding of the dummy variable lyst is the
list (a b c); but as soon as the answer c is passed along to the function or person that
requested the rac, the value of the atom lyst used in getting that answer is totally
forgotten. The Lisp interpreter will tell you that lyst is an "unbound atom" if you ask
for its value. Applicative programmers much prefer lambda bindings to ordinary setq
bindings.
I personally am not a fanatic about avoiding setq's and other functions that cause side
effects. Though I find the applicative style to be jus-telegant, I find it impractical
when it comes to the construction of large Al-style programs. Therefore I shall not
advocate the applicative style here, though I shall adhere to it when possible. Strictly
speaking, in applicative programming, you cannot even define new functions, since a
def statement causes a permanent change to take place in the genie's memory-namely,
the permanent storage in memory of the function definition. So the ideal applicative
approach would have functions, like variable bindings, being created only
temporarily, and their definitions would be discarded the moment after they had been
used. But this is extreme "applicativism".
         For your edification, here are a few more simple function definitions.

-> (def rdc (lambda (lyst) (reverse (cdr (reverse lyst)))))
-> (def snot (lambda (x lyst) (reverse (cons x (reverse lyst)))))
-> (def twice (lambda (n) (plus n n)))

The functions rdc and snoc are analogous to cdr and cons, only backwards. Thus, the
rdc of (a b c d e) is (a b c d), and if you type (snoc 5 '(1 2 3 4)), you will get (1 2 3 4
5) as your answer.

                                        *   *   *




Lisp: Atoms and Lists                                                                  406

All, of this is mildly interesting so far, but if you want to see the genie do anything
truly surprising, you have to allow it to make some decisions based on things that
happen along the way. These are sometimes called “conditional wishes". A typical
example would be the following:

-> (mod ((eq x 1) 'land) ((eq x 2) 'sea))

The value returned by this statement will be the atom land if x has value 1, and the
atom sea if x has value 2. Otherwise, the value returned will be nil (i.e., if x is 5). The
atom eq (pronounced "eek") is the name of ..a common Lisp function that returns the
atom t (standing for "true") if its two arguments have the same value, and nil (for "no"
or "false") if they do not.
        A cond statement is a list whose car is the function name cond, followed by
any number of cond clauses, each of which is a two-element list. The first element of
each clause is called its condition, the second element its result. The clauses'
conditions are checked out by the Lisp genie one by one, in order; as soon as it finds a
clause whose condition is "true" (meaning that the condition returns anything other
than nil!), it begins calculating that clause's result, whose value gets returned as the
value of the whole cond statement. None of the further clauses is even so much as
glanced at! This may sound more complex than it ought to. The real idea is no more
complex than saying that it looks for the first condition that is satisfied, then it returns
the corresponding result.
        Often one wants to have a catch-all clause at the end whose condition is sure
to be satisfied, so that, if all other conditions fail, at least this one will be true and the
accompanying result, rather than nil, will be returned. It is easy as pie to make a
condition whose value is non-nil; just choose it to be t for instance, as in the
following:

       (Cond ((eq x 1) 'land)
       ((eq x 2) 'sea)
       (t 'air))


'Depending on what the value of x is, we will get either land, sea, or air as the value
of this cond, but we'll never get nil. Now here are a few sample cond statements for
you to play genie to:

-> (cond ((eq (oval pi) pie) (oval (snot pie pi)))
       (t (eval (snoc (rac pi) pi))))
-> (cond ((eq 2 2) 2) ((eq 3 3) 3))
-> (cond (nil 'no-no-no)
       ((eq '(car nil) '(cdr nil)) 'hmmm)
       (t 'yes-yes-yes))




Lisp: Atoms and Lists                                                                    407

The answers are: 8, 2, and yes-yes-yes. Did you notice that (car nil) and (cdr nil)
were quoted?
I shall close this portion of the column by displaying a patterned family of function
definitions, so obvious in their pattern that you would think that the Lisp genie would
just sort of "get the hang of it" after seeing the first few .... Unfortunately, though,
Lisp genies are frustratingly dense (or at least they play at being dense), and they will
not jump to any conclusion unless it has been completely spelled out. Look first at the
family:

-> (def square (lambda (k) (times k k)))
-> (def cube (lambda (k) (times k (square k))))
-> (def 4th-power (lambda (k) (times k (cube k))))
-> (def 5th-power (lambda (k) (times k (4th-power k))))
-> (def 6th-power (lambda (k) (times k (5th-power k))))
-> .
-> .
-> .
-> .

My question for you is this: Can you invent a definition for a two parameter function
that subsumes all of these in one fell swoop? More concretely, the question is: How
would one go about defining a two-parameter function called power such that, for
instance, (power 9 3) yields 729 on being evaluated, and (power 7 4) yields 2,401 ? I
have supplied you, in this column, with all the necessary tools to do this, provided you
exercise some ingenuity.

                                        *   *   *

I thought I would end this column with a newsbreak about a freshly discovered beast-
the homely Glazunkian porpuquine, so called because it is found only on the island of
Glazunkia (claimed by Upper Bitbo, though it is just off the coast of Burronymede).
And what is a porpuquine, you ask? Why, it's a strange breed of porcupine, whose
quills-of which, for some reason, there are always exactly nine (in Outer Glazunkia)
or seven (in Inner Glazunkia)-are smaller porpuquines. Oho! This would certainly
seem to be an infinite regress! But no. It's just that I forgot to mention that there is a
smallest size of porpuquine: the zero-inch type, which, amazingly enough, is totally
bald of quills. So, quite luckily (or perhaps unluckily, depending on your point of
view), that puts a stop to the threatened infinite regress. This remarkable beast is
shown in a rare photograph in Figure 17-1.
        Students of zoology might be interested to learn that the quills on 5-inch
porpuquines are always 4-inch porpuquines, and so on down the line. And students of
anthropology might be equally intrigued to know that the residents of Glazunkia (both
Outer and Inner) utilize the nose (yes, the nose) of the zero-inch porpuquine as a unit
of barter-an odd thing to our




Lisp: Atoms and Lists                                                                 408

FIGURE 17-1. The homely Inner Glazunkian porpuquine, Porpuquinus
verdimontianus. The size of this particular specimen has not been ascertained,
although it appears to be at least a 4-incher. The buying power of a porpuquine is the
number of zero-inch noses on it. Larger noses, oddly enough, are worth nothing.
[Photograph by David J. Moser. ]

minds; but then, who are you and I to question the ancient wisdom of the Outer and
Inner Glazunkians? Thus, since a largish porpuquine-say a 3-incher or 4-incher-
contains many, many such tiny noses, it is a most valuable commodity. The value of a
porpuquine is sometimes referred to as its "buying power", or just "power" for short.
For instance, a 2-incher found in Inner Glazunkia is almost twice as powerful as a 2-
incher found in Outer Glazunkia. Or did I get it backward? It's rather confusing!
        Anyway, why am I telling you all this? Oh, I just thought you'd like to hear
about it. Besides, who knows? You just might wind up visiting Glazunkia .(Inner or
Outer) one of these fine days. And then all of this could come in mighty handy.




Lisp: Atoms and Lists                                                             409

